% Source: https://artofproblemsolving.com/community/c3h2586363_number_theory
% Stable identity: aops-topic-2586363-q01
For $n=1$ the value is $5+12+12+3=32$, which does not work.

For $n\ge2$, $5^n=25\cdot5^{n-2}$, and $5^{n-2}\equiv1\pmod4$, so $5^n\equiv25\pmod{100}$. The condition becomes
\[12n^2+12n+28\equiv0\pmod{100}\iff4\left(3n^2+3n+7\right)\equiv0\pmod{100}\iff3n^2+3n+7\equiv0\pmod{25}.\]
\emph{Necessary:} modulo $5$, $3n^2+3n+2$ takes the values $2,3,0,3,2$ for $n\equiv0,1,2,3,4$, so $n\equiv2\pmod5$.

\emph{Sufficient:} if $n=5t+2$, then
\[3n^2+3n+7=75t^2+60t+12+15t+6+7=75t^2+75t+25\equiv0\pmod{25}.\]
So the answer is $\boxed{\text{all }n\equiv2\pmod5}$, i.e. $n=2,7,12,17,\dots$. (For example $n=2$: $25+48+24+3=100$.)
